\section{Instruction Data Landscape: quién escribe el prompt y quién escribe la completion}

Una vez establecido el objetivo de entrenamiento, el siguiente paso es decidir de dónde provienen las demonstrations. En clase no se dividieron los datos simplemente en «buenos/malos»; en cambio, se compararon el costo, la escalabilidad, la carga de filtrado y la herencia de estilo a lo largo de un espectro continuo que va de model-generated a human-written.

\subsection{Estructura básica de una muestra}

La sección anterior ya distinguió tres tipos de señales de supervisión; esta sección primero descompone una demonstration de SFT para responder en qué puntos exactos la fuente de los datos entra en el modelo. Al leer las dos diapositivas siguientes, conviene rastrear por separado quién produce el prompt y quién produce la completion, en lugar de evaluar solo si la respuesta final es fluida. La unidad mínima de SFT suele ser instruction/context más completion. Esta definición parece sencilla, pero oculta dos decisiones fundamentales: si el prompt proviene de la distribución real de usuarios y si la completion la escribe una persona, un modelo fuerte o un modelo débil por autoarranque; estas dos fuentes determinan la cobertura de capacidades y la transmisión de sesgos.

\lecturefigure{slide-07.jpg}{Estructura instance/completion de una instruction demonstration}{Diapositiva oficial p.7}

\begin{knowledgebox}{Lectura de la figura: primero se separan input provenance y output provenance}
En la figura, el bloque azul es instruction/context y el bloque verde es la demonstration. Ambos pueden provenir de productores distintos: por ejemplo, una persona escribe el prompt y GPT-4 escribe la answer, o un modelo genera ambas partes. Al juzgar la calidad de los datos no basta con leer la fluidez de la answer; también hay que verificar si el prompt cubre las tareas de los usuarios objetivo.
\end{knowledgebox}

\lecturefigure{slide-08.jpg}{Espectro continuo de los datos de SFT, de model-generated a human-written}{Diapositiva oficial p.8}

\begin{knowledgebox}{Lectura de la figura: el eje horizontal no es una simple clasificación por calidad}
Los métodos de la izquierda son baratos y escalables, pero dependen más de la seed, del teacher model y del filtrado; los de la derecha son costosos, pero pueden incorporar intenciones humanas reales y una expresión concisa. El eje horizontal muestra la authorship: no garantiza que estar más a la derecha sea mejor ni que la synthetic data sea necesariamente de baja calidad.
\end{knowledgebox}

\teachervoice{Rajani insiste en que la diferencia entre Surge, OpenAssistant, Dolly y Self-Instruct es, ante todo, «quién escribe el input y quién escribe el output». No está elaborando un dataset leaderboard, sino construyendo una provenance taxonomy.}

\subsection{Tres mecanismos de Synthetic Data}

Aunque todos pertenecen a la model-generated data, el human role no es el mismo en Self-Instruct, UltraChat y CAMEL. Al compararlos lado a lado se entiende que las variables centrales de un synthetic pipeline son la seed, retrieval/material, los agent roles y el filtering, y no un simple interruptor de «generar con GPT».

\lecturefigure{slide-09.jpg}{Self-Instruct: seed tasks, bootstrapping y filtering}{Diapositiva oficial p.9}

\begin{knowledgebox}{Lectura de la figura: la cadena de expansión y la cadena de errores son la misma}
Unas pocas seeds escritas por personas se amplían mediante few-shot generation y, tras pasar por task classification y filtrado, entran en el conjunto de tareas. La capacidad de expansión proviene del teacher model; el colapso de patrones, las repeticiones y los errores también se propagan desde ese mismo teacher, por lo que el filtering no es un paso de limpieza opcional, sino parte del algoritmo de generación.
\end{knowledgebox}

\lecturefigure{slide-10.jpg}{UltraChat: human-in-the-loop que selecciona materiales y genera de forma iterativa}{Diapositiva oficial p.10}

\begin{knowledgebox}{Lectura de la figura: las personas no necesariamente escriben la respuesta final, pero también pueden cambiar la distribución de los datos}
Las personas buscan temas y materiales de apoyo, y deciden las tareas y la dirección de las preguntas de seguimiento; un modelo fuerte genera el contenido de los diálogos. En comparación con el autoarranque puramente basado en modelos, este mecanismo sitúa el trabajo humano en la selección de temas y el refinamiento iterativo, y a cambio de un menor costo de redacción obtiene un control más explícito de la cobertura.
\end{knowledgebox}

\lecturefigure{slide-11.jpg}{CAMEL: role-playing de user/assistant para generar diálogos de varios turnos}{Diapositiva oficial p.11}

\begin{knowledgebox}{Lectura de la figura: la interacción entre dos agents sigue restringida por la tarea inicial}
Un modelo desempeña el papel de AI assistant y otro el de user, y ambos desarrollan la conversación por su cuenta en torno a una task de alto nivel. Es adecuado para generar interaction traces a gran escala, pero el prompt de cada rol, la termination y los sesgos compartidos por los modelos determinan la diversidad de los diálogos; no se debe equiparar automáticamente «multi-agent» con «usuarios más reales».
\end{knowledgebox}

\teachervoice{En clase, los tres métodos se ordenaron como líneas de producción con distinto «grado de participación humana»: Self-Instruct se apoya en seed/filter, UltraChat hace que las personas aporten materiales y refinen repetidamente, y CAMEL solo pide a las personas definir la tarea de alto nivel. Con esto, el ponente muestra que la synthetic data no es un tipo de datos, sino una familia de workflows.}

\subsection{Del Landscape a la elección de H4}

En ese momento, los datos públicos crecían con rapidez, pero el equipo aun así decidió comprar instruction demonstrations completamente escritas por personas. Esta decisión no niega la synthetic data; su propósito era obtener una línea base humana controlable para después comparar, en experimentos posteriores, el task mix, la length y la quantity.

\lecturefigure{slide-12.jpg}{Versión actualizada del instruction dataset landscape de 2023}{Diapositiva oficial p.12}

\begin{knowledgebox}{Lectura de la figura: los nombres de los conjuntos de datos requieren una explicación de su mecanismo}
Self-Instruct/Unnatural Instructions se basan principalmente en la expansión con modelos; T0/FLAN convierten NLP tasks existentes al instruction format; UltraChat/ShareGPT dependen de diálogos de modelos o de usuarios; Dolly/OpenAssistant/Surge están más cerca de la redacción humana. El conocimiento verdaderamente transferible es el mecanismo y la provenance, no memorizar nombres.
\end{knowledgebox}

\lecturefigure{slide-13.jpg}{H4 elige Surge-Instruct dentro del SFT data continuum}{Diapositiva oficial p.13}

\begin{importantbox}{Por qué comprar todavía datos humanos}
El equipo quería usar «human-written» como grupo de control experimental: si el estilo, los errores factuales o las preferencias de longitud de un synthetic teacher ya están en los datos, es difícil determinar si la mejora del modelo final proviene de la cobertura de tareas o de la teacher imitation. Los datos humanos son costosos, pero ofrecen una estructura de errores distinta.
\end{importantbox}

\teachervoice{El ponente dice sin rodeos que humans inefficient and expensive, pero que tampoco se debe subestimar la calidad de la manually created data. Que H4 eligiera a principios de 2023 una ruta «muy manual» fue un diseño de investigación sujeto a restricciones de presupuesto, no una recomendación permanente para todos los proyectos.}

\begin{knowledgebox}{Términos clave: fuentes de datos comunes}
\begin{tabularx}{\textwidth}{@{}L{0.18\textwidth}L{0.25\textwidth}X@{}}
\toprule
Nombre & Problema que resuelve & Posición del mecanismo en esta clase \\
\midrule
Self-Instruct & Ampliar un instruction set a partir de pocas seeds & Generación con modelo + filtrado por clasificación \\
UltraChat & Diálogos de varios turnos a gran escala & Personas eligen temas/materiales + generación con modelo fuerte \\
CAMEL & Trayectorias de interacción con roles & Dos LLM en role-playing \\
FLAN/T0 & Unificar el instruction format multitarea & Reescritura de tareas y mixture \\
LIMA & Pocas demostraciones de alta calidad & human-written, cantidad pequeña \\
OpenAssistant y Dolly & Diálogos humanos de la comunidad o de empleados & human-authored baseline \\
Surge-Instruct & Demonstrations humanas controlables de H4 & vendor workforce + distribución especificada \\
\bottomrule
\end{tabularx}
\end{knowledgebox}

\subsection{Resumen de la sección}

La pregunta fundamental de la instruction data es la authorship y el workflow: quién decide la tarea, quién produce el contenido y quién filtra. La synthetic data y la human data conllevan costos y sesgos distintos, y los experimentos posteriores deben registrar esta provenance de forma explícita.
